% corpo.tex — conteúdo comum dos MWEs web e print (mwe_book_web.tex,
% mwe_book_print.tex). Do início da Introdução à lista de referências; capa,
% folha de rosto, sumário e colofão ficam nos .tex principais.
% Texto sem U+0060 (a Atkinson não tem o glifo) e sem hifenização manual.

\chapter{Introdução}

Este documento exercita a \emph{linha book} acessível (versão 0.6.0): corpo em
Atkinson Hyperlegible Next, sem hifenização, sem órfãs nem viúvas, alinhamento à
esquerda e pé de página irregular. O Código de Ética \parencite{cfp2005} e a
literatura sobre eventos acessíveis \parencite{exemplo2020,exemplo2024} servem de
pretexto para chamadas de citação no texto: na opção \texttt{print} elas aparecem
entre parênteses; na opção \texttt{web} a lista de referências traz todas as
entradas, citadas ou não.

\iniciotextual

\chapter{Deficiência intelectual, deficiência psicossocial e transtorno do espectro autista: orientações para a realização de eventos e apresentações acessíveis}

Este capítulo tem título longo de propósito: ele é a entrada mais comprida do
sumário e exercita a margem direita das linhas de capítulo, que sem hifenização
não pode contar com a quebra de palavra para fugir do número da página.

\section{Preparação do ambiente}

Escolha um espaço com boa acústica, iluminação uniforme e sinalização clara.
Informe com antecedência o roteiro da programação, os intervalos e as pausas. O
cabeçalho do documento acompanha o capítulo corrente, e a numeração árabe das
páginas começa neste capítulo.

Palavras longas, como sociodemográfico, inconstitucionalíssimo, otorrinolaringologista
e psicopedagogicamente, ficam inteiras: a linha termina antes delas. Um traço de
união como em ético-profissional ou socioeducativo também não é ponto de quebra.

\section{Durante a apresentação}

Fale de frente para o público, com ritmo moderado, e descreva oralmente o que
estiver nos slides. Evite siglas sem explicação e frases longas demais. Uma
pessoa com deficiência intelectual lê melhor quando cada parágrafo traz uma só
ideia, e uma pessoa que usa leitor de tela depende da ordem de leitura correta
dos títulos, listas e tabelas.

Mais um parágrafo comprido para forçar quebras de página e testar as regras de
órfãs e viúvas: nenhuma primeira linha de parágrafo fica sozinha no pé da página
e nenhuma última linha fica sozinha no topo da seguinte. Repetimos a ideia com
outras palavras para dar volume: a acessibilidade não é um acréscimo no fim do
processo, é uma decisão de projeto desde o início, que envolve o espaço, o
conteúdo, a comunicação e as pessoas que participam do evento, de ponta a ponta.

\chapter{Estrangeirismos e termos técnicos}

Termos de outra língua levam o idioma no \emph{Span} da árvore de tags, para o
leitor de tela trocar de voz. Exemplos: o formato \estrangeirismo{en}{Easy Read}
para textos em leitura facilitada, a \estrangeirismo{en}{quiet room} para
pessoas que precisam de um espaço silencioso e a pausa \estrangeirismo{en}{coffee break}
programada entre as falas.

Um termo já incorporado ao português, como \emph{slide}, leva só o itálico.

\chapter{Tabela de dados}

A tabela abaixo usa o ambiente \texttt{tabelaguia}: a primeira linha é
cabeçalho (repetido quando a tabela quebra de página) e as linhas alternam o
fundo por \verb|\rowcolor|.

\begin{tabelaguia}{L{0.26\linewidth}L{0.28\linewidth}L{0.30\linewidth}}
\rowcolor{crp1!15}\textbf{Barreira} & \textbf{Exemplo} & \textbf{Providência}\\\endhead
Comunicacional & Falas rápidas, sem descrição das imagens & Descrever oralmente o conteúdo visual\\
Atitudinal & Tratar a pessoa com infantilização & Dirigir-se à própria pessoa\\
\rowcolor{crp1!6}Arquitetônica & Degraus na entrada & Rampa e piso tátil\\
Digital & Documento sem estrutura de títulos & Marcar títulos e ordem de leitura\\
Metodológica & Atividade com tempo único para todas as pessoas & Prever tempos e formatos alternativos\\
\rowcolor{crp1!6}Programática & Regra do evento que impede acompanhante & Rever o regulamento\\
Instrumental & Material só em papel & Oferecer formato digital acessível\\
\rowcolor{crp1!6}Sensorial & Música alta durante o intervalo & Oferecer espaço silencioso\\
Comunicacional & Vídeo sem legenda & Legenda e janela de Libras\\
Atitudinal & Pressa para concluir a fala da pessoa & Respeitar o tempo de cada fala\\
\rowcolor{crp1!6}Digital & Formulário sem rótulos nos campos & Rotular cada campo\\
Arquitetônica & Banheiro estreito & Adequar à norma de acessibilidade\\
Metodológica & Avaliação só escrita & Admitir a resposta oral ou gravada\\
Programática & Inscrição sem opção de adaptação & Perguntar as necessidades no cadastro\\
\end{tabelaguia}

\chapter{Texto corrido}

Este capítulo existe para dar volume e provocar quebras de página em lugares
variados. A psicologia, como ciência e profissão, fundamenta-se no respeito à
diversidade e na valorização da singularidade de cada pessoa. Por isso as
práticas de comunicação e de divulgação do conhecimento devem refletir esses
valores, garantindo que todas as pessoas, sem exceção, possam participar
plenamente dos espaços de diálogo e de aprendizado promovidos pelo Conselho.

Adotar práticas acessíveis não é só uma questão de conformidade legal, mas um
ato de respeito e de responsabilidade social. Ao eliminar as barreiras
comunicacionais, atitudinais, digitais, metodológicas, arquitetônicas, físicas e
programáticas, contribuímos para um ambiente profissional mais ético,
acolhedor e verdadeiramente inclusivo, e para a universalização do acesso da
população ao conhecimento da ciência psicológica.

A leitura em tela e a leitura em papel pedem cuidados distintos. Na tela, o
documento rola ou passa de página em página, sem encadernação, e a margem externa
larga, que no livro impresso é a margem do polegar, aqui é só uma faixa de
respiro. No papel, a página par e a página ímpar se espelham, e o capítulo abre
sempre à direita, mesmo que isso deixe uma página em branco antes dele.

Uma sequência de parágrafos curtos ajuda a observar o pé irregular da página.
Primeiro parágrafo curto, de duas linhas, para testar a regra de que um
parágrafo de duas linhas passa inteiro para a página seguinte quando não cabe.

Segundo parágrafo curto, também de duas linhas, com palavras suficientes para
quebrar em duas linhas exatas na largura da mancha em A5 de onze pontos.

Terceiro parágrafo curto, igualmente de duas linhas, para fechar a sequência e
não deixar a página terminar exatamente no fim de um parágrafo.

\chapter{Considerações finais}

Capítulo curto: no \texttt{print}, a página seguinte é aberta à direita, e se
este capítulo terminar numa página ímpar sobra uma página em branco antes da
lista de referências.
